\section{Задание}
Вычислить приближённо интеграл по формуле Ньютона-Котеса с данным значением $n$, оценить погрешность приближения с помощью остаточного члена.

Вычислить приближённо интеграл по формуле Гаусса с данным значением $n$, оценить погрешность приближения по более точному значению интеграла, вычисленному по формуле Ньютона-Лейбница.

\medskip
\noindent\textbf{Вариант 47:}
\[
    f(x)=x\sin x, \quad x\in[0;\pi].
\]

Для формулы Гаусса:
\[
    n=5.
\]

Для формулы Ньютона-Котеса:
\[
    n=4.
\]

\section{Цель работы}
Научиться приближённо оценивать значения определённых интегралов по квадратурным формулам Ньютона-Котеса и Гаусса.

\section{Ход работы}

Требуется вычислить определённый интеграл
\[
I=\int\limits_0^{\pi}x\sin x\,dx.
\]

\subsection{Теоретические сведения}

Задача численного интегрирования состоит в приближённом вычислении определённого интеграла
\[
I=\int\limits_a^b f(x)\,dx.
\]

Квадратурная формула заменяет интеграл конечной суммой вида
\[
\int\limits_a^b f(x)\,dx \approx \sum_{i=0}^{n-1} A_i f(x_i),
\]
где $x_i$~--- узлы квадратурной формулы, а $A_i$~--- её веса.

\subsubsection*{Формула Ньютона-Котеса}

Формулы Ньютона-Котеса являются интерполяционными квадратурными формулами. Подынтегральная функция заменяется интерполяционным многочленом, после чего вычисляется интеграл от этого многочлена.

В данной работе для формулы Ньютона-Котеса задано $n=4$, то есть используется формула трёх восьмых. Пусть
\[
h=\frac{b-a}{3}, \quad x_i=a+ih, \quad i=0,1,2,3.
\]

Тогда формула имеет вид
\[
I_{3/8}=\frac{3h}{8}\left(f(x_0)+3f(x_1)+3f(x_2)+f(x_3)\right).
\]

Остаточный член формулы трёх восьмых оценивается следующим образом:
\[
|R|\leq \frac{3}{80}h^5M_4,
\]
где
\[
M_4=\max_{x\in[a;b]}|f^{(4)}(x)|.
\]

\subsubsection*{Формула Гаусса}

Для вычисления интеграла по формуле Гаусса с $n$ узлами на отрезке $[-1;1]$ используется формула
\[
\int\limits_{-1}^{1} f(t)\,dt \approx \sum_{i=0}^{n-1}w_i f(t_i),
\]
где $t_i$~--- узлы формулы Гаусса, а $w_i$~--- соответствующие веса.

Для перехода к произвольному отрезку $[a;b]$ используется замена
\[
x_i=\frac{a+b}{2}+\frac{b-a}{2}t_i.
\]

Тогда
\[
\int\limits_a^b f(x)\,dx \approx \frac{b-a}{2}\sum_{i=0}^{n-1}w_i f(x_i).
\]

\subsection{Точное значение интеграла}

Найдём точное значение интеграла по формуле Ньютона-Лейбница:
\[
I=\int\limits_0^{\pi}x\sin x\,dx.
\]

Воспользуемся интегрированием по частям:
\[
u=x, \quad dv=\sin x\,dx,
\]
\[
du=dx, \quad v=-\cos x.
\]

Тогда
\[
\int x\sin x\,dx=-x\cos x+\int\cos x\,dx=-x\cos x+\sin x.
\]

Следовательно,
\[
I=\left(-x\cos x+\sin x\right)\Big|_0^{\pi}.
\]

Подставим пределы интегрирования:
\[
I=(-\pi\cos\pi+\sin\pi)-(-0\cdot\cos0+\sin0)=\pi.
\]

Таким образом,
\[
I=\pi\approx 3.141593.
\]

\subsection{Вычисление интеграла по формуле Ньютона-Котеса}

Для формулы Ньютона-Котеса задано $n=4$, поэтому используем формулу трёх восьмых.

Имеем
\[
a=0, \quad b=\pi.
\]

Шаг равен
\[
h=\frac{b-a}{3}=\frac{\pi-0}{3}=\frac{\pi}{3}\approx 1.047198.
\]

Узлы квадратурной формулы:
\[
x_0=0, \quad x_1=\frac{\pi}{3}, \quad x_2=\frac{2\pi}{3}, \quad x_3=\pi.
\]

Вычислим значения функции в узлах:
\[
f(x)=x\sin x.
\]

\[
\begin{aligned}
&f(x_0)=f(0)=0\cdot\sin0=0,\\
&f(x_1)=f\left(\frac{\pi}{3}\right)=\frac{\pi}{3}\sin\frac{\pi}{3}=\frac{\pi\sqrt3}{6}\approx0.906900,\\
&f(x_2)=f\left(\frac{2\pi}{3}\right)=\frac{2\pi}{3}\sin\frac{2\pi}{3}=\frac{\pi\sqrt3}{3}\approx1.813799,\\
&f(x_3)=f(\pi)=\pi\sin\pi=0.
\end{aligned}
\]

Подставим полученные значения в формулу трёх восьмых:
\[
I_{3/8}=\frac{3h}{8}\left(f(x_0)+3f(x_1)+3f(x_2)+f(x_3)\right).
\]

\[
\begin{aligned}
I_{3/8}={}&\frac{3\cdot\frac{\pi}{3}}{8}
\left(0+3\cdot0.906900+3\cdot1.813799+0\right)=\\
={}&\frac{\pi}{8}(0+2.720699+5.441398+0)=\\
={}&3.205248.
\end{aligned}
\]

Следовательно,
\[
I_{3/8}\approx3.205248.
\]

\subsection{Оценка погрешности формулы Ньютона-Котеса}

Фактическая погрешность равна
\[
\Delta I_{3/8}=|I-I_{3/8}|.
\]

Подставим значения:
\[
\Delta I_{3/8}=|3.141593-3.205248|=0.063655.
\]

Оценим погрешность с помощью остаточного члена. Для этого найдём четвёртую производную функции.

\[
f(x)=x\sin x.
\]

\[
f'(x)=\sin x+x\cos x,
\]
\[
f''(x)=2\cos x-x\sin x,
\]
\[
f'''(x)=-3\sin x-x\cos x,
\]
\[
f^{(4)}(x)=x\sin x-4\cos x.
\]

На отрезке $[0;\pi]$ имеем
\[
M_4=\max_{x\in[0;\pi]}|x\sin x-4\cos x|\approx4.7773.
\]

Тогда
\[
|R|\leq\frac{3}{80}h^5M_4.
\]

Подставим значения:
\[
|R|\leq\frac{3}{80}\left(\frac{\pi}{3}\right)^5\cdot4.7773\approx0.225607.
\]

Получаем
\[
\Delta I_{3/8}=0.063655<0.225607.
\]

Следовательно, оценка погрешности выполнена.

На листинге \ref{lst:newton_cotes} приведена программа, вычисляющая интеграл по формуле Ньютона-Котеса.

\begin{listing}[H]
\caption{Вычисление интеграла по формуле Ньютона-Котеса}
\vspace{-20pt}
\label{lst:newton_cotes}
\begin{minted}[linenos, numbersep=10pt, firstnumber=1]{matlab}
f = @(x) x .* sin(x);

a = 0;
b = pi;

h = (b - a) / 3;
x = [a, a + h, a + 2 * h, b];
y = f(x);

I_nc = 3 * h / 8 * (y(1) + 3 * y(2) + 3 * y(3) + y(4));
I_exact = pi;
error_nc = abs(I_exact - I_nc);

fprintf('I_nc = %.6f\n', I_nc);
fprintf('error_nc = %.6f\n', error_nc);
\end{minted}
\end{listing}

\subsection{Вычисление интеграла по формуле Гаусса}

Для формулы Гаусса задано $n=5$. Используем узлы и веса формулы Гаусса на отрезке $[-1;1]$.

\begin{table}[H]
    \centering
    \caption{Узлы и веса формулы Гаусса при $n=5$}
    \label{table:gauss_nodes}
    \begin{tblr}{
        colspec = {X[c]X[c]X[c]},
        hlines,
        vlines,
        rows = {abovesep=5pt,belowsep=5pt, $},
        row{1} = {font=\normalfont}
    }
        i & t_i & w_i \\
        0 & -0.906180 & 0.236927 \\
        1 & -0.538469 & 0.478629 \\
        2 & 0.000000 & 0.568889 \\
        3 & 0.538469 & 0.478629 \\
        4 & 0.906180 & 0.236927 \\
    \end{tblr}
\end{table}

Пересчитаем узлы на отрезок $[0;\pi]$:
\[
x_i=\frac{0+\pi}{2}+\frac{\pi-0}{2}t_i=\frac{\pi}{2}+\frac{\pi}{2}t_i.
\]

Получим:
\[
\begin{aligned}
&x_0=0.147372, \quad x_1=0.724971, \quad x_2=1.570796,\\
&x_3=2.416622, \quad x_4=2.994220.
\end{aligned}
\]

Вычислим значения функции и произведения $w_if(x_i)$.

\begin{table}[H]
    \centering
    \caption{Вычисление слагаемых формулы Гаусса}
    \label{table:gauss_calc}
    \begin{tblr}{
        colspec = {X[c]X[c]X[c]X[c]},
        hlines,
        vlines,
        rows = {abovesep=5pt,belowsep=5pt, $},
        row{1} = {font=\normalfont}
    }
        i & x_i & f(x_i) & w_if(x_i) \\
        0 & 0.147372 & 0.021640 & 0.005127 \\
        1 & 0.724971 & 0.480738 & 0.230095 \\
        2 & 1.570796 & 1.570796 & 0.893609 \\
        3 & 2.416622 & 1.602495 & 0.767000 \\
        4 & 2.994220 & 0.439670 & 0.104170 \\
    \end{tblr}
\end{table}

Формула Гаусса на отрезке $[0;\pi]$ имеет вид
\[
I_G=\frac{\pi}{2}\sum_{i=0}^{4}w_if(x_i).
\]

Подставим найденные значения:
\[
\begin{aligned}
I_G={}&\frac{\pi}{2}(0.005127+0.230095+0.893609+0.767000+0.104170)=\\
={}&\frac{\pi}{2}\cdot2.000000\approx3.141593.
\end{aligned}
\]

Более точное значение, полученное вычислением в программе:
\[
I_G=3.141593.
\]

На листинге \ref{lst:gauss} приведена программа, вычисляющая интеграл по формуле Гаусса.

\begin{listing}[H]
\caption{Вычисление интеграла по формуле Гаусса}
\vspace{-20pt}
\label{lst:gauss}
\begin{minted}[linenos, numbersep=10pt, firstnumber=1]{matlab}
f = @(x) x .* sin(x);

a = 0;
b = pi;

nodes = [-0.9061798459, -0.5384693101, 0, ...
          0.5384693101,  0.9061798459];
weights = [0.2369268851, 0.4786286705, 0.5688888889, ...
           0.4786286705, 0.2369268851];

x = (a + b) / 2 + (b - a) / 2 * nodes;
y = f(x);

I_gauss = (b - a) / 2 * sum(weights .* y);
I_exact = pi;
error_gauss = abs(I_exact - I_gauss);

fprintf('I_gauss = %.12f\n', I_gauss);
fprintf('error_gauss = %.12f\n', error_gauss);
\end{minted}
\end{listing}

\subsection{Оценка погрешности формулы Гаусса}

Погрешность формулы Гаусса оценим по точному значению интеграла:
\[
\Delta I_G=|I-I_G|.
\]

Так как
\[
I=\pi\approx3.141592653590,
\]
а
\[
I_G\approx3.141592826824,
\]
то
\[
\Delta I_G=|3.141592653590-3.141592826824|\approx0.000000173.
\]

Следовательно, формула Гаусса при $n=5$ дала очень точное приближение.

\section{Вывод}

В результате выполнения лабораторной работы был вычислен интеграл
\[
\int\limits_0^{\pi}x\sin x\,dx.
\]

Точное значение интеграла равно
\[
I=\pi\approx3.141593.
\]

По формуле Ньютона-Котеса при $n=4$ было получено значение
\[
I_{3/8}=3.205248,
\]
при этом фактическая погрешность составила
\[
\Delta I_{3/8}=0.063655.
\]

По формуле Гаусса при $n=5$ было получено значение
\[
I_G=3.141593,
\]
при этом погрешность составила
\[
\Delta I_G\approx0.000000173.
\]

Таким образом, формула Гаусса в данном случае дала более точное приближение значения интеграла.
